\documentclass[UTF8,zihao=-4]{ctexart}
\usepackage[a4paper,margin=1.85cm]{geometry}
\usepackage{hyperref}
\usepackage{verbatim}
\usepackage{xcolor}
\usepackage{fancyhdr}
\hypersetup{hidelinks}
\setlength{\parindent}{0pt}
\setlength{\parskip}{0.45em}
\emergencystretch=4em
\sloppy
\pagestyle{fancy}
\fancyhf{}
\lhead{期末复习 PPT 专项模拟卷}
\rhead{\thepage}
\title{补充模拟卷 02：语义分析、AST、符号表补缺}
\author{}
\date{}
\begin{document}
\maketitle
\setcounter{tocdepth}{1}
\tableofcontents
\newpage


\section*{A 卷题目}
\addcontentsline{toc}{section}{A 卷题目}

\subsection*{一、依赖图与拓扑求值顺序（20 分）}
\addcontentsline{toc}{subsection}{一、依赖图与拓扑求值顺序（20 分）}

考虑语义规则：

\begin{quote}
\footnotesize
\begin{verbatim}
E -> E1 + T
  E.val = E1.val + T.val

E -> T
  E.val = T.val

T -> num
  T.val = num.lexeme
\end{verbatim}
\end{quote}

对输入：

\begin{quote}
\footnotesize
\begin{verbatim}
3 + 4
\end{verbatim}
\end{quote}

\noindent 1. 写出主要属性依赖关系。\\
\noindent 2. 写出一种合法的属性求值顺序。\\
\noindent 3. 计算最终 \texttt{E.val}。\\

\subsection*{二、受控副作用（20 分）}
\addcontentsline{toc}{subsection}{二、受控副作用（20 分）}

考虑声明翻译：

\begin{quote}
\footnotesize
\begin{verbatim}
D -> T id
T -> int
\end{verbatim}
\end{quote}

语义动作：

\begin{quote}
\footnotesize
\begin{verbatim}
addType(id.name, T.type)
\end{verbatim}
\end{quote}

\noindent 1. 说明该语义动作的副作用是什么。\\
\noindent 2. 为什么该动作必须在 \texttt{T.type} 已经确定之后执行。\\
\noindent 3. 写出受控副作用应满足的基本要求。\\

\subsection*{三、CST 与 AST（20 分）}
\addcontentsline{toc}{subsection}{三、CST 与 AST（20 分）}

表达式：

\begin{quote}
\footnotesize
\begin{verbatim}
a + b * c
\end{verbatim}
\end{quote}

\noindent 1. 说明 CST 和 AST 的区别。\\
\noindent 2. 画出该表达式的 AST。\\
\noindent 3. 说明为什么后续语义分析更常使用 AST。\\

\subsection*{四、符号表与作用域（20 分）}
\addcontentsline{toc}{subsection}{四、符号表与作用域（20 分）}

考虑程序片段：

\begin{quote}
\footnotesize
\begin{verbatim}
{
  int x;
  {
    float x;
  }
}
\end{verbatim}
\end{quote}

\noindent 1. 说明进入和退出代码块时符号表栈如何变化。\\
\noindent 2. 在内层代码块中查找 \texttt{x} 时应得到哪个声明。\\
\noindent 3. 退出内层代码块后查找 \texttt{x} 时应得到哪个声明。\\

\subsection*{五、类型信息与标识符绑定（20 分）}
\addcontentsline{toc}{subsection}{五、类型信息与标识符绑定（20 分）}

回答：

\noindent 1. 符号表通常记录哪些信息。\\
\noindent 2. AST 中的标识符结点为什么要绑定符号表条目。\\
\noindent 3. 符号表对类型检查有什么作用。\\

\section*{B 简明答案}
\addcontentsline{toc}{section}{B 简明答案}

\subsection*{一、依赖图与拓扑求值顺序答案}
\addcontentsline{toc}{subsection}{一、依赖图与拓扑求值顺序答案}

对 \texttt{3 + 4}，主要属性依赖：

\begin{quote}
\footnotesize
\begin{verbatim}
T1.val 依赖 num(3).lexeme
E1.val 依赖 T1.val
T2.val 依赖 num(4).lexeme
E.val  依赖 E1.val 和 T2.val
\end{verbatim}
\end{quote}

一种合法求值顺序：

\begin{quote}
\footnotesize
\begin{verbatim}
1. T1.val = 3
2. E1.val = T1.val = 3
3. T2.val = 4
4. E.val = E1.val + T2.val = 3 + 4 = 7
\end{verbatim}
\end{quote}

最终：

\begin{quote}
\footnotesize
\begin{verbatim}
E.val = 7
\end{verbatim}
\end{quote}

\subsection*{二、受控副作用答案}
\addcontentsline{toc}{subsection}{二、受控副作用答案}

副作用：

\begin{quote}
\footnotesize
\begin{verbatim}
addType(id.name, T.type) 会修改符号表。
\end{verbatim}
\end{quote}

必须在 \texttt{T.type} 确定后执行，因为符号表中需要记录标识符的正确类型。如果 \texttt{T.type} 尚未计算，就无法把正确类型写入符号表。

受控副作用要求：

\begin{quote}
\footnotesize
\begin{verbatim}
1. 执行位置明确。
2. 执行顺序确定。
3. 不破坏属性依赖关系。
4. 不引入依赖环。
5. 不导致重复插入或错误类型绑定。
\end{verbatim}
\end{quote}

\subsection*{三、CST 与 AST 答案}
\addcontentsline{toc}{subsection}{三、CST 与 AST 答案}

CST：

\begin{quote}
\footnotesize
\begin{verbatim}
具体语法树，完整反映文法推导过程，包含括号、辅助非终结符等语法细节。
\end{verbatim}
\end{quote}

AST：

\begin{quote}
\footnotesize
\begin{verbatim}
抽象语法树，去掉冗余语法细节，保留表达式和语句的核心结构。
\end{verbatim}
\end{quote}

\texttt{a + b * c} 的 AST：

\begin{quote}
\footnotesize
\begin{verbatim}
      +
     / \
    a   *
       / \
      b   c
\end{verbatim}
\end{quote}

后续语义分析更常使用 AST，因为 AST 更直接表示程序语义结构，便于类型检查、符号绑定、优化和代码生成。

\subsection*{四、符号表与作用域答案}
\addcontentsline{toc}{subsection}{四、符号表与作用域答案}

进入外层代码块：

\begin{quote}
\footnotesize
\begin{verbatim}
push 外层符号表
插入 x : int
\end{verbatim}
\end{quote}

进入内层代码块：

\begin{quote}
\footnotesize
\begin{verbatim}
push 内层符号表
插入 x : float
\end{verbatim}
\end{quote}

内层代码块中查找 \texttt{x}：

\begin{quote}
\footnotesize
\begin{verbatim}
先查内层符号表，得到 x : float
\end{verbatim}
\end{quote}

退出内层代码块：

\begin{quote}
\footnotesize
\begin{verbatim}
pop 内层符号表
\end{verbatim}
\end{quote}

此后查找 \texttt{x}：

\begin{quote}
\footnotesize
\begin{verbatim}
得到外层的 x : int
\end{verbatim}
\end{quote}

退出外层代码块：

\begin{quote}
\footnotesize
\begin{verbatim}
pop 外层符号表
\end{verbatim}
\end{quote}

\subsection*{五、类型信息与标识符绑定答案}
\addcontentsline{toc}{subsection}{五、类型信息与标识符绑定答案}

符号表通常记录：

\begin{quote}
\footnotesize
\begin{verbatim}
名字
类型
作用域
存储位置
参数信息
函数信息
\end{verbatim}
\end{quote}

AST 中标识符结点绑定符号表条目的原因：

\begin{quote}
\footnotesize
\begin{verbatim}
同名标识符可能出现在不同作用域中。
绑定后可以明确该标识符对应哪一个声明。
\end{verbatim}
\end{quote}

符号表对类型检查的作用：

\begin{quote}
\footnotesize
\begin{verbatim}
通过符号表查出变量、函数、表达式的类型，
判断赋值、运算、函数调用等是否类型匹配。
\end{verbatim}
\end{quote}
\end{document}
